\chapter{Dicas e Boas Práticas} \label{chp:ap_dicas}

Esse apêndice tem como objetivo explicar como construir um texto científico, apresentando dicas na construção do texto e demonstrando exemplos de códigos em LaTeX para as principais necessidades na elaboração do texto.

% ---
\section{Inclusão de outros arquivos}
% ---

É uma boa prática dividir o seu documento em diversos arquivos, 
e não
apenas escrever tudo em um único. 
Esse recurso foi utilizado neste documento. 
Para incluir diferentes arquivos em um arquivo principal, 
de modo que cada arquivo incluído fique em uma página diferente, utilize o comando:

\begin{verbatim}
   \include{documento-a-ser-incluido}      % sem a extensão .tex
\end{verbatim}

Para incluir documentos sem quebra de páginas, utilize:

\begin{verbatim}
   \input{documento-a-ser-incluido}      % sem a extensão .tex
\end{verbatim}

% ---
\section{Citações}
% ---

AQUI FALAR SOBRE CITACOES

% ---
\section{Imagens e Tabelas}\label{sec:LABEL_CHP_2_SEC_A}
% ---

Toda tabela~\ref{tab:MATRIX_NOTAS_EXEMPLO}

\begin{table}[!htb]
  \centering
  \caption{Fragmento de uma matriz de notas de um sistema de recomendação de filmes.}
  \label{tab:MATRIX_NOTAS_EXEMPLO}
  \rowcolors{1}{}{lightlightgray}
  \begin{tabular}{l c c c}
  \toprule
    & Titanic & Poderoso Chefão & Matrix \\
    \midrule
        Filipe Braida & 4 & $\emptyset$ & 3  \\
        Leandro Alvim & 4 & 5 & 5  \\
        Bruno Dembogurski & 4 & 5 & 5  \\
        Fellipe Duarte & $\emptyset$ & 5 & $\emptyset$  \\ 
    \bottomrule
  \end{tabular}
\end{table}

\section{Images}\label{sec:LABEL_CHP_2_SEC_B}
Reference: \url{http://en.wikibooks.org/wiki/LaTeX/Importing_Graphics}

\begin{figure}
  \centering
  \includegraphics[width=0.6\textwidth]{imagens/chick.png}
  \caption{Chick}
  \label{fig:LABEL_FIG_1}
\end{figure}

\begin{algorithm}[H]
\floatname{algorithm}{Algoritmo}

\textbf{Entrada} conjunto de notas $R$, limiar $\tau$, modelo de previsão $\varphi$.

\textbf{Saída} conjunto de notas sem ruídos $R^{*}$.

$R^{*} \gets \{\}$

\For{$(u, i, r) \in R$}{
$\tilde{r} \gets \varphi(u,i)$\;

\If{$|\tilde{r} - r| < \tau$}{
$R^{*} \gets R^{*} \cup \{(u,i,r)\}$
}
}
\caption{Filtragem das avaliações com ruído proposto por \cite{OMahony2006}.}
\label{alg:mahony}
\end{algorithm}

\section{Equações}
Reference: \url{http://en.wikibooks.org/wiki/LaTeX/Mathematics}

Also: \url{http://en.wikibooks.org/wiki/LaTeX/Advanced_Mathematics}

\begin{equation}
  (x + y)^2 = x^2 + 2xy + y^2
  \label{eq:LABEL_EQ_1}
\end{equation}

\section{Listings}\label{sec:LABEL_CHP_2_SEC_D}
Reference: \url{http://en.wikibooks.org/wiki/LaTeX/Source_Code_Listings}

\codec{C}{alg:LABEL_CODE_1}{codigos/codigo-c.txt}

\codejava{Java}{alg:LABEL_CODE_2}{codigos/codigo-java.txt}

\section{References}\label{sec:LABEL_CHP_2_SEC_E}
\begin{itemize}
  \item Referencing \refchapter{chp:LABEL_CHP_1}
  \item Referencing \refsection{sec:LABEL_CHP_1_SEC_A}
  \item Referencing \refsection{sec:LABEL_CHP_1_SEC_C}
  \item Referencing \reftable{tab:LABEL_TAB_1}
  \item Referencing \reffigure{fig:LABEL_FIG_1}
  \item Referencing \refequation{eq:LABEL_EQ_1}
  \item Referencing \reflisting{alg:LABEL_CODE_1}
  \item Article \cite{braida2015transforming}
  \item Referencing \refappendix{chp:LABEL_APP_1}
\end{itemize}